\documentclass[12pt,a4paper,UTF8]{ctexart}
\usepackage[a4paper,top=22mm,bottom=22mm,left=22mm,right=22mm,headheight=15pt]{geometry}
\usepackage{amsmath,amssymb,mathtools,bm}
\usepackage{graphicx}
\usepackage{xcolor}
\usepackage{booktabs,longtable,array}
\usepackage{enumitem}
\usepackage{fancyhdr}
\usepackage{caption}
\usepackage{float}
\usepackage{microtype}
\usepackage{setspace}
\usepackage[hidelinks,unicode]{hyperref}
\usepackage[most]{tcolorbox}

\definecolor{ContestBlue}{HTML}{17365D}
\definecolor{SoftBlue}{HTML}{EEF4FA}
\definecolor{SoftGray}{HTML}{F5F6F7}
\hypersetup{pdftitle={第99套},pdfauthor={物理竞赛教学资料}}

\setstretch{1.24}
\setlength{\parindent}{2em}
\setlength{\parskip}{0.35em}
\setlist{itemsep=0.35em,topsep=0.4em}
\captionsetup{font=small,labelfont=bf,justification=centering}
\setcounter{secnumdepth}{0}

\pagestyle{fancy}
\fancyhf{}
\fancyhead[L]{\small 第99套}
\fancyfoot[C]{\thepage}
\renewcommand{\headrulewidth}{0.4pt}

\ctexset{
  section={format=\Large\bfseries\color{ContestBlue},beforeskip=1.2em,afterskip=0.7em},
  subsection={format=\large\bfseries\color{ContestBlue},beforeskip=1.0em,afterskip=0.5em},
  subsubsection={format=\normalsize\bfseries,beforeskip=0.8em,afterskip=0.35em}
}

\makeatletter
\newcommand{\tightlist}{\setlength{\itemsep}{0pt}\setlength{\parskip}{0pt}}
\newcommand{\maxwidth}{\ifdim\Gin@nat@width>\linewidth\linewidth\else\Gin@nat@width\fi}
\newcommand{\maxheight}{\ifdim\Gin@nat@height>0.72\textheight 0.72\textheight\else\Gin@nat@height\fi}
\newcommand{\pandocbounded}[1]{\begingroup\setkeys{Gin}{width=\maxwidth,height=\maxheight,keepaspectratio}#1\endgroup}
\makeatother

\begin{document}

\begin{titlepage}
  \centering
  \vspace*{58mm}
  {\Huge\bfseries\color{ContestBlue} 第99套\par}
  \vfill
  \vspace*{28mm}
\end{titlepage}

\section{第 1 题}\label{ux7b2c-1-ux9898}

在一台显示压力的台秤上放置一个装水的容器。水的密度为 \(\rho_1\)、质量为
\(m_1\)，容器质量为 \(m_2\)。将质量为 \(m_3\)、密度为 \(\rho_3\)
的物体完全浸没在水中并由静止释放，其中 \(\rho_3>\rho_1\)，物体随后下沉。

\begin{enumerate}
\def\labelenumi{\arabic{enumi}.}
\tightlist
\item
  求物体下沉过程中，容器、水和物体所组成系统的质心加速度大小。
\item
  求物体下沉过程中台秤的示数。
\item
  将物体换成质量为 \(m_4\)、密度为 \(\rho_4\) 的均匀链条，其中
  \(\rho_4>\rho_1\)。链条最上端贴近水面，完全浸没后由静止释放。水深与链长均为
  \(h\)，故释放瞬间链条最下端恰好接触容器底。链条与容器底的碰撞为完全非弹性碰撞，求全过程中台秤示数的最大值。
\end{enumerate}

忽略水流速度引起的压强变化及水流对液面形状的影响。

\begin{figure}[H]
\centering
\includegraphics[width=0.9\linewidth,height=\textheight,keepaspectratio,alt={图 1-1　物体在水中下沉，以及均匀链条逐段落至容器底部。}]{../figures/question-01.pdf}
\caption*{图 1-1　物体在水中下沉，以及均匀链条逐段落至容器底部。}
\end{figure}

\newpage

\section{第 2 题}\label{ux7b2c-2-ux9898}

一个质量为 \(M\) 的均匀实心圆锥，轴线与母线的夹角为 \(\beta\)，母线长为
\(l\)；圆锥轴线与竖直方向的夹角记为 \(\theta\)。

\begin{enumerate}
\def\labelenumi{\arabic{enumi}.}
\tightlist
\item
  求圆锥质心到圆锥尖端的距离。
\item
  圆锥尖端朝下立于水平地面，受微小扰动后从竖直状态开始倾倒。尖端与地面间的静摩擦足够大，使尖端始终不滑动。取
  \(\beta=\pi/4\)，求保证倾倒过程中不滑动所需的最小静摩擦因数。
\item
  若圆锥初始绕自身轴线以很大的角速度 \(\omega\)
  自转，轴线与竖直方向夹角为 \(\theta\)，求轴线的稳定进动角速度
  \(\Omega\)。
\end{enumerate}

\begin{figure}[H]
\centering
\includegraphics[width=0.9\linewidth,height=\textheight,keepaspectratio,alt={图 2-1　以尖端为支点的圆锥及轴线、母线和竖直方向之间的角度。}]{../figures/question-02.pdf}
\caption*{图 2-1　以尖端为支点的圆锥及轴线、母线和竖直方向之间的角度。}
\end{figure}

\newpage

\section{第 3 题}\label{ux7b2c-3-ux9898}

电路如图所示。已知

\[R_1=3R,\qquad R_2=R,\qquad C_1=C_2=C,\qquad L_1=L_2=L,\]

电源电动势为
\(\mathcal E\)。开关起初闭合，电路已达到直流稳态，电压表视为理想电压表。

\begin{enumerate}
\def\labelenumi{\arabic{enumi}.}
\tightlist
\item
  求稳态时电压表的示数 \(U_1\)。
\item
  突然断开开关，求断开后一瞬间电压表的示数 \(U_2\)。
\item
  忽略电压表及其他元件中的能量损失，求从断开开关到电路达到新稳态的过程中，电阻
  \(R_1\)、\(R_2\) 上分别产生的焦耳热 \(Q_1\)、\(Q_2\)。
\end{enumerate}

\begin{figure}[H]
\centering
\includegraphics[width=0.684\linewidth,height=\textheight,keepaspectratio,alt={图 3-1　电路结构：含两电阻、两电容、两电感及理想电压表的桥式电路。}]{../figures/question-03.pdf}
\caption*{图
3-1　电路结构：含两电阻、两电容、两电感及理想电压表的桥式电路。}
\end{figure}

\newpage

\section{第 4 题}\label{ux7b2c-4-ux9898}

一种同轴电缆由两层金属组成：内导体是半径为 \(r_1\)
的无限长实心圆柱，外导体是内、外半径分别为 \(r_2\)、\(r_3\)
的无限长同轴圆柱壳；两导体之间为真空。总电流 \(I\)
沿内导体流向纸面外，并沿外导体反向返回。

\begin{enumerate}
\def\labelenumi{\arabic{enumi}.}
\tightlist
\item
  若内、外导体中的电流密度均匀，求磁感应强度随轴线距离 \(r\) 的分布
  \(B(r)\)。
\item
  若导体理想，电流只分布在金属表面：内导体电流位于 \(r=r_1\)
  的表面，外导体回流位于 \(r=r_2\) 的内表面，求 \(B(r)\)。
\item
  在第（2）问条件下，求电缆单位长度内储存的磁场能量。
\item
  仍按第（2）问，内导体固定不动。外圆柱壳因材料原因缓慢发生径向变化，并保持密度
  \(\rho\) 不变。若内半径由 \(r_2\) 变为 \(r_2+x\)，其中
  \(|x|\ll r_2\)，求外半径的相应变化。
\item
  在题设电磁作用远大于金属弹性作用的近似下，判断第（4）问的径向运动是否为简谐振动；若是，求周期
  \(T\)；若不是，说明原因。
\end{enumerate}

可用小量展开

\[\ln(1+x)=x-\frac{x^2}{2}+O(x^3).\]

\begin{figure}[H]
\centering
\includegraphics[width=0.53\linewidth,height=\textheight,keepaspectratio,alt={图 4-1　同轴电缆横截面及三个特征半径 r\_1,r\_2,r\_3。}]{../figures/question-04.pdf}
\caption*{图 4-1　同轴电缆横截面及三个特征半径 \(r_1,r_2,r_3\)。}
\end{figure}

\newpage

\section{第 5 题}\label{ux7b2c-5-ux9898}

两根平行、光滑且电阻不计的导轨间距为
\(l\)，空间中存在垂直纸面向里的匀强磁场 \(\bm B\)。两根质量均为
\(m\)、电阻不计的导体棒垂直跨接在导轨上，棒间连接一根劲度系数为 \(k\)
的轻弹簧。初始时弹簧恰为原长，两棒均静止。记左右两棒向右的速度分别为
\(v_1\)、\(v_2\)，并在初始时刻给左棒一个向右的冲量 \(J\)。

\begin{enumerate}
\def\labelenumi{\arabic{enumi}.}
\tightlist
\item
  忽略回路中的电感与电容，求 \(v_1(t)\)。注意初始时刻可能出现冲击电流。
\item
  若回路自感恒为 \(L\)，求 \(v_1(t)\)。
\item
  若回路自感恒为 \(L\)，并在回路中串联电容 \(C\)，求
  \(v_1(t)\)。为简化结果，设
  \[\frac1{LC}=\frac{k}{m},\qquad \frac{B^2l^2}{m}=\frac1C,\] 并用
  \(k/m\) 表示答案中的特征频率。
\end{enumerate}

\begin{figure}[H]
\centering
\includegraphics[width=0.9\linewidth,height=\textheight,keepaspectratio,alt={图 5-1　匀强磁场中的双导体棒、平行导轨和连接弹簧。}]{../figures/question-05.pdf}
\caption*{图 5-1　匀强磁场中的双导体棒、平行导轨和连接弹簧。}
\end{figure}

\newpage

\section{第 6 题}\label{ux7b2c-6-ux9898}

一个长方形绝热容器被分为 \(A\)、\(B\)、\(C\)、\(D\) 四个腔室。\(A\)
室左壁导热并与温度为 \(T_0\) 的大热源接触，其余外壁绝热。\(A\)、\(B\)
间为自由绝热活塞；\(B\)、\(C\)
间为固定绝热隔板，其中嵌有只允许氦气通过、不允许氮气通过的半透膜；\(C\)、\(D\)
间为自由绝热活塞。

\(A\) 室有 \(1\,\mathrm{mol}\) 氦气，\(B\)、\(C\) 两室合计有
\(2\,\mathrm{mol}\) 氦气，\(C\) 室另有 \(1\,\mathrm{mol}\) 氮气，\(D\)
室有 \(1\,\mathrm{mol}\) 氮气。初始时四室体积均为
\(V_0\)。气体均视为理想气体，氦气为单原子气体，氮气为双原子气体，过程准静态。

\begin{enumerate}
\def\labelenumi{\arabic{enumi}.}
\tightlist
\item
  求初始平衡时 \(A\)、\(B\)、\(C\)、\(D\) 四室的温度之比。
\item
  缓慢给 \(D\) 室加热，直至其压强变为初始压强的 \(2\)
  倍，求新的四室温度之比。
\item
  求第（2）问中向 \(D\) 室输入的热量。
\end{enumerate}

\begin{figure}[H]
\centering
\includegraphics[width=0.9\linewidth,height=\textheight,keepaspectratio,alt={图 6-1　装置结构：四个腔室、两自由绝热活塞、固定半透膜及恒温热源。}]{../figures/question-06.pdf}
\caption*{图
6-1　装置结构：四个腔室、两自由绝热活塞、固定半透膜及恒温热源。}
\end{figure}

\newpage

\section{第 7 题}\label{ux7b2c-7-ux9898}

一个内壁光滑的圆环形气缸，管道截面半径为 \(a\)，中心圆半径为
\(r\)。气缸内等间距放置 \(N\) 个绝热活塞，其中 \(N\) 很大；活塞质量均为
\(m\)，与内壁无摩擦且不漏气。已知

\[r\gg \frac rN\gg a,\]

因此每个小气室可近似看作直圆柱。初始平衡时，相邻活塞之间均有
\(\nu\,\mathrm{mol}\) 氮气，温度均为 \(T_0\)。第 \(i\)
个活塞相对平衡位置的位移为 \(x_i\)，且 \(|x_i|\ll r/N\)。

\begin{enumerate}
\def\labelenumi{\arabic{enumi}.}
\tightlist
\item
  在线性近似下，求第 1、2 个活塞间气体的压强变化量。
\item
  考虑所有活塞均以圆频率 \(\omega\)
  作简谐运动、相位随序号线性变化的基频波模式。相邻活塞相位差
  \(\Delta\phi\ll1\)，可用
  \(\cos\Delta\phi\simeq1-(\Delta\phi)^2/2\)。并假设所得波速远小于气体声速，使每个小气室内压强始终近似均匀。求该模式的最小相速度。
\end{enumerate}

\begin{figure}[H]
\centering
\includegraphics[width=0.53\linewidth,height=\textheight,keepaspectratio,alt={图 7-1　圆环形气缸内的等间距绝热活塞。}]{../figures/question-07.pdf}
\caption*{图 7-1　圆环形气缸内的等间距绝热活塞。}
\end{figure}

\newpage

\section{第 8 题}\label{ux7b2c-8-ux9898}

在一次高能质子碰撞中产生顶夸克 \(t\) 和反顶夸克 \(\bar t\)。其中

\[t\to b+W^+,\qquad W^+\to\mu^++\nu.\]

除顶夸克和 \(W^\pm\)
外，忽略其余粒子的静止质量。中微子的横向动量可由横向动量守恒确定，但其
\(z\) 分量无法直接测量。某次 ATLAS 事例中，与 \(W^+\)
衰变有关的测量数据如下，动量单位均为 \(\mathrm{GeV}/c\)：

\begin{center}
\begin{tabular}{lrrr}
\toprule
粒子 & $p_x$ & $p_y$ & $p_z$ \\
\midrule
反缪子 $\mu^+$ & $-24.7$ & $-24.9$ & $-12.4$ \\
Jet 1（$b$ 喷注） & $-14.2$ & $50.1$ & $94.1$ \\
中微子 $\nu$ & $-104.1$ & $5.3$ & 未知 \\
\bottomrule
\end{tabular}
\end{center}

已知 \(W^+\) 玻色子的静质量为

\[m_W=80.4\,\mathrm{GeV}/c^2,\]

求中微子纵向动量 \(p_z^{(\nu)}\) 的所有可能值，以 \(\mathrm{GeV}/c\)
表示。

\begin{figure}[H]
\centering
\includegraphics[width=0.516\linewidth,height=\textheight,keepaspectratio,alt={图 8-1　已知粒子的横向动量；中微子纵向分量不由此图确定。}]{../figures/question-08.pdf}
\caption*{图 8-1　已知粒子的横向动量；中微子纵向分量不由此图确定。}
\end{figure}

\end{document}
